\begin{textsample}{NEG Dim 7 – global\_north\_2025\_03 – Score -1.00 – global\_north\_2025\_03/122409360.txt}  \label{ex:f7_neg_139}
Palestinians mourn over the bodies of their relatives killed by an Israeli airstrike in Beit Lahiya, northern Gaza Strip on Saturday.
The Israeli military says it was targeting"terrorists"operating a drone in the area.
But the head of Gaza ' s civil defense, Mahmoud Basal, says they were aid workers from the London-based Al Khair Foundation, along with two journalists working with the group.
Jehad Alshrafi/AP hide caption toggle caption Jehad Alshrafi/AP TEL AVIV, Israel -- Two back-to-back Israeli airstrikes in the northern Gaza city of Beit Lahiya killed at least nine people on Saturday, according to Gaza civil defense.
Witnesses and Gaza health officials say several of those killed were aid workers and journalists.
The Israeli military says it was targeting two terrorists operating a drone in the area.
In a statement, it says it hit those individuals, and then struck again when"a number of additional terrorists"collected the drone operating equipment and entered a vehicle.
But the head of Gaza civil defense @ @ @ @ @ @ @ @ @ @ from the London-based Al-Khair Foundation, along with several journalists working with the group."They had been wanting to build some tents for the displaced people in the area, and they had been using the drone to shoot and take coordinates of the place,"Basal says.
Sponsor Message A local worker with the Al-Khair Foundation confirmed those details to NPR, saying the workers were taking aid to the neighborhood for a Ramadan Iftar celebration.
The worker spoke on the condition of anonymity as they were not authorized to speak to the media.
The area where the strike happened is designated as a ' free movement area ' by the Israeli military, far from the"buffer zone"along the outer edges of Gaza where movement is restricted.
The Israeli military did not immediately respond to NPR ' s request for further comment on the incident.
In a statement on the militant ' s group ' s Telegram channel, Hamas called the strikes an"escalation"and a"deliberate sabotage of any @ @ @ @ @ @ @ @ @ @ The attack comes as negotiations for the continuation of the fragile ceasefire between Israel and Hamas, which began on January 19, have stalled.
The first phase of the deal expired earlier this month, and Hamas has been pushing for the second phase to begin.
Israel is refusing, and calling for a new plan, which does not have a provision for a long-term end to the conflict -- something Hamas is adamant about.
This is not the first time that Palestinians have been killed by Israeli fire during the nearly \textbf{two-month} ceasefire, but it is the deadliest single incident thus far.
More than 150 Palestinians have been killed by Israel during the ceasefire, according to the Gaza Ministry of Health.
Sponsor Message Meanwhile, Israel ' s blockade of aid on Gaza -- including electricity -- is entering its third week.
Israel imposed the blockade as a way to pressure Hamas to agree to the deal that it wants.
International aid groups say the blockade is having devastating effects on the two million Palestinians @ @ @ @ @ @ @ @ @ @ than 15 months of war, calling it collective punishment and a violation of international law.

% matched lemmas: two-month
\end{textsample}
